
En
effet, soient @@M0@@ le point g\'en\'erique de @@M1@@ et @@M2@@ la cl\^oture alg\'ebrique de @@M3@@. Il est clair que @@M4@@ est la limite du syst\`eme projectif filtrant @@M5@@ des sch\'emas affines int\`egres munis d'un morphisme fini et fid\`element plat @@M6@@ o\`u @@M7@@ est un ouvert non-vide de @@M8@@. Le normalis\'e de @@M9@@ provient donc d'un @@M10@@-sch\'ema @@M11@@ pour un @@M12@@, le morphisme fini et surjectif @@M13@@ d'un @@M14@@-morphisme fini et surjectif @@M15@@ (EGA IV 8.8.2, 8.10.5). Enfin, comme @@M16@@ est normal et int\`egre, il existe un ouvert @@M17@@ de @@M18@@ tel que @@M19@@ soit normal et int\`egre sur @@M20@@ (EGA IV 9.9.5).

\textbf{Remarque 3.10.} Conservant les notations et conditions de 3.9, supposons qu'il existe un morphisme r\'esolvant @@M21@@ (EGA IV 7.9.1), (par exemple si @@M22@@, d'apr\`es Abhyankar). Alors un raisonnement similaire montre qu'il existe un morphisme fini et fid\`element plat @@M23@@, un sch\'ema @@M24@@ lisse sur @@M25@@ et un morphisme propre surjectif birationnel @@M26@@.

Dans la d\'emonstration de la premi\`ere partie de 3.8, on s'est ramen\'e au cas o\`u @@M27@@ est int\`egre et @@M28@@ est normal et int\`egre sur @@M29@@, donc @@M30@@ est int\`egre. Si @@M31@@, alors @@M32@@ et l'assertion est triviale. Sinon, quitte \`a restreindre @@M33@@, on peut supposer que pour tout @@M34@@, on a @@M35@@, donc @@M36@@ est @@M37@@-r\'egulier et la dimension @@M38@@ de @@M39@@ est @@M40@@.

Soit alors @@M41@@ une famille limit\'ee de classes de faisceaux inversibles sur les fibres de @@M42@@. Soit @@M43@@ un faisceau inversible sur une fibre g\'eom\'etrique @@M44@@ de @@M45@@, avec le polyn\^ome de Hilbert @@M46@@, alors le polyn\^ome de Hilbert de @@M47@@ est @@M48@@. Lorsque la classe de @@M49@@ parcourt @@M50@@, c-\`a-d la
